% Proposed supplement paragraph, to follow \subsection{Envelope-mean subharmonicity}
% (supp:envelope-proof). Klaus, 2026-11-24, from Kaiti's envelope_boundary_scope.md.
% Claim concerns interior MAXIMA. One closing sentence on minima.

\paragraph{Interior maxima of the excursion and the beat amplitude.}
The product of the polarizations obeys
$\nabla^2(\varepsilon_1\varepsilon_2)=2\nabla\varepsilon_1\cdot\nabla\varepsilon_2$,
whose sign is unrestricted, and an excursion, a temporal maximum minus a temporal
minimum, need not be subharmonic. A local example shows that both can have a strict
three-dimensional interior maximum under the conditions above. In dimensionless
coordinates, take $0<c<1/2$, $q=x^2-\tfrac12(y^2+z^2)$, and the harmonic potentials
\[
\phi_{1,2}=-x\mp\tfrac12(x^2-y^2)-c\left[\tfrac13x^3-\tfrac12x(y^2+z^2)\right].
\]
Their fields along the fixed $x$ direction are
$\varepsilon_{1,2}=-\partial_x\phi_{1,2}=1\pm x+cq$, positive near the origin. Then
\[
\varepsilon_1\varepsilon_2=1+(2c-1)x^2-c(y^2+z^2)+O(|\mathbf x|^4),
\]
with a negative-definite Hessian at the origin, and the envelope excursion
$2\min(\varepsilon_1,\varepsilon_2)=2(1+cq-|x|)$ is also strictly maximal there.
This is an existence example, not a prediction for a montage. The subharmonic maps
above can still have interior minima; the restriction concerns maxima only.
